\section{Blockierfunktion}
Dieser Abschnitt widmet sich der Blockierfunktion. Zunächst wird der Funktionshintergrund und die Funktionsweise dieser Funktion beschrieben. Im Anschluss erfolgt eine detaillierte Analyse der spezifischen Anforderungen sowie der bestehenden Softwareimplementierung (C-Code). Auf Grundlage der gewonnenen Erkenntnisse wird die Schnittstelle für den Playground erweitert und die Funktion in Matlab modelliert.

\subsection{Funktionsbeschreibung}
Das Blockieren eines Werkzeugs kann zu erhöhten Strömen in den Motorphasen führen, was eine übermäßige Erwärmung der Motorwicklungen zur Folge hat. Um einer Überlastung vorzubeugen, ermöglicht die Blockierfunktion eine Reduzierung des Anlaufstroms bei übermäßiger Belastung. Dies wird durch einen dynamischen Blockierzähler realisiert, der bei jeder erkannten Blockade inkrementiert wird. Im Stillstand oder während des Betriebs wird der Zähler nach einer vordefinierten Zeit schrittweise dekrementiert. Die Aktivierung oder Deaktivierung der Stromreduzierung erfolgt basierend auf dem Wert des Blockierzählers. Zusätzlich schützt eine einstellbare Abschaltzeit das Gerät vor Überhitzung, indem es bei anhaltendem Blockieren des Werkzeugs abgeschaltet wird \cite{SW_ZUSE}. Ein beispielhafter zeitlicher Verlauf dieser Blockierfunktion ist in \autoref{png_funktion_blk} dargestellt.
\begin{figure}[H]
	\centering
	\includegraphics[width=0.97\textwidth]{./Bilder/Funktionsbeschreibung_Blk.png}
	\caption{Beispielhafter Zeitablauf der Blockierfunktion \cite{Lastenheft}}
	\label{png_funktion_blk}
\end{figure}
Der Blockierzähler, dargestellt im unteren Diagramm, wird bei jeder erkannten Blockade schrittweise erhöht, bis sein maximaler Wert (im gezeigten Beispiel neun) erreicht ist. Eine schrittweise Verringerung des Zählers erfolgt nach einem definierten Zeitintervall, das abhängig vom Betriebszustand des Motors (laufend oder stillstehend) variiert. Die Aktivierung der Stromreduktion erfolgt, sobald der Zähler einen festgelegten Schwellenwert (im gezeigten Beispiel sieben) überschreitet und wird deaktiviert, sobald der Zähler den unteren Schwellenwert (im Beispiel null) erreicht.

\subsection{Analyse des Lastenhefts und des C-Codes}
\label{subsec_blk_analyse}
In diesem Abschnitt werden die spezifischen Anforderungen aus dem Lastenheft bei STIHL für die Blockierfunktion aufgeführt. Diese Anforderungen bilden die Grundlage für die Implementierung des entsprechenden C-Codes in der Firmware. Durch eine ergänzende Analyse des C-Codes kann die Modellierung der Funktion vereinfacht und gleichzeitig sichergestellt werden, dass deren Funktionalität vollständig abgebildet wird. \\
Die nachfolgenden Anforderung sind wörtlich aus dem Master-Lastenheft für Akkugeräte entnommen \cite{Lastenheft}:
\begin{itemize}
\item Der Schutz vor Überhitzung durch Blockieren wird durch einen dynamischen Blockierzähler realisiert. Der Zähler hat einen Wertebereich von 0 bis \textit{CpsBlkWrtBlockadeMax}.
\item Wenn ein Werkzeug länger als \textit{CpsBlkZeiBlockierungErkannt} nach der Erkennung blockiert wird, erhöht das Gerät den Blockierzähler um 1.
\item Während der Motor läuft reduziert das Gerät den Zähler nach Ablauf der Zeit \textit{CpsBlkWrtBlockadeDekrementBetrieb} um 1. 
\item Während der Motor steht reduziert das Gerät den Zähler nach Ablauf der Zeit \textit{CpsBlkWrtBlockadeDekrementStillstand} um 1.
\item Wenn der Blockierzähler beim Inkrementieren den Wert \textit{CpsBlkWrtBlockadeStromReduktionEin} erreicht, begrenzt das Gerät den Strom auf den Wert \textit{CpsBlkStrMotorReduzierterAnlauf}.
\item Wenn der Blockierzähler beim Dekrementieren den Wert \textit{CpsBlkWrtBlockadeStromReduktionAus} erreicht, setzt das Gerät die Strombegrenzung zurück.
\item Falls die Zeit \textit{CpsBlkZeiBlockierungAbschalten} nach einer erkannten Blockierung abgelaufen ist, soll das Gerät den Fehler \textit{MC\_ERR\_BLOCKIERERKENNUNG} ausgeben. Der Fehler wird zurückgesetzt, sobald ein erneuter Startversuch erfolgt („Kundenwunsch Gerätestart“ - abhängig von Bedienlogik) und dabei der Fehler nicht weiter anliegt. Die Überprüfung auf das Anliegen des Fehlers wird im Rahmen des erneuten Startversuchs überprüft und falls er weiter vorliegt, erneut geschrieben.
\end{itemize}
Der C-Code in der Firmware basiert auf den genannten Anforderungen. Eine zusätzliche Untersuchung des C-Codes ermöglicht es, sowohl die Informationen zum Ein- und Ausschalten der Funktion über maschinenabhängige Optionen als auch die zusätzlichen Eingangssignale, die für eine optimale Funktionalität der Funktion erforderlich sind, zu identifizieren. \\
Aus der Analyse des C-Codes \cite{Firmware_gesamt} ergeben sich folgende Punkte für die Modellierung der Blockierfunktion:
\begin{itemize}
\item Die Funktion gliedert sich in die Blockadeerkennung und die Blockadeauswertung. Die Blockade kann durch zwei unabhängige Erkennungsverfahren identifiziert werden. Ein Verfahren stellt eine sicherheitsrelevante Erkennung dar, die innerhalb der Basissoftware durchgeführt wird und dort verbleiben soll. Das andere Verfahren variiert je nach eingesetztem Motortyp (DC oder EC). Da die meisten modernen Akkugeräte von STIHL mit EC-Motoren ausgestattet sind, wird in dieser Arbeit ausschließlich die Blockadeerkennung für EC-Motoren in Matlab modelliert. Diese Erkennung kann über die maschinenabhängige Option \lstinline{ACTIVATE_BLOCKADE_MONITORING} im Header-File aktiviert oder deaktiviert werden.

\item Die Blockadeerkennung für EC-Motoren erfolgt durch die Abfrage des elektrischen Winkels des Motordrehfeldes. Wird innerhalb einer festgelegten Zeitspanne \textit{CpsBlkZeiBlockierungErkannt} der Delta-Wert des Drehfelds von \unit[60]{$^\circ$} nicht überschritten, so wird eine Blockade detektiert. Die Abfrage des Winkels erfolgt mit einer Abtastzeit von \unit[62.5]{µs}, um eine präzise Erkennung zu gewährleisten. Die Auswertung der Daten erfolgt jedoch zyklisch in einem Intervall von \unit[1]{ms}.

\item Die Unterscheidung zwischen langsamem und schnellem Dekrementieren des Zählers basiert auf der Abfrage der aktuellen Motordrehzahl im Vergleich zur Start-Motordrehzahl. Die Ausgangs-Motordrehzahl stellt dabei die minimale Drehzahl der derzeitigen Maschinenkonfiguration dar. Überschreitet die aktuelle Drehzahl der Maschine die Startdrehzahl, kommt die Zeitspanne \textit{CpsBlkWrtBlockadeDekrementBetrieb} zum Einsatz, um das Dekrementieren vorzunehmen. Andernfalls wird die Zeitspanne \textit{CpsBlkWrtBlockadeDekrementStillstand} für das Dekrementieren verwendet.

\item Der Blockierzähler kann während des Betriebs der Maschine nur einmal inkrementiert werden. Für eine erneute Inkrementierung ist ein kurzfristiger Motorstillstand erforderlich. Diese Bedingung wird in der Basissoftware durch eine Abfrage des Motorsteuerstatus realisiert.

\end{itemize}

\subsection{Erweiterung der Schnittstelle für den Playground}
In diesem Abschnitt wird die Erweiterung der Schnittstelle zur korrekten Übergabe der Ein- und Ausgangssignale zwischen der Basissoftware und dem Autocode behandelt. Zu diesem Zweck wird die Datei \textit{Playground.c} angepasst. Die erforderlichen Eingangssignale werden aus der Basissoftware in der Datei \textit{Playground.c} aufgerufen und anschließend an den Playground übergeben. Der entsprechende C-Code wird in dieser Arbeit nicht dargestellt. Stattdessen ist in \autoref{png_io_blk} ein schematisches Diagramm für die Blockierfunktion mit allen Ein- und Ausgängen abgebildet.
\begin{figure}[H]
	\centering
	\includegraphics[width=1\textwidth]{./Bilder/IO_Blk.pdf}
	\caption{Ein- und Ausgänge der Blockierfunktion}
	\label{png_io_blk}
\end{figure}
Ein Teil der Blockadeerkennung bei EC-Motoren wird mit einer Abtastzeit von \unit[62.5]{µs} ausgeführt. Da der Playground jedoch nur einmal pro Millisekunde aufgerufen wird, muss dieser Teil außerhalb des Playgrounds verarbeitet werden. Dafür dient der Eingang \textit{Neuer Sektor detektiert}, der anzeigt, ob der elektrische Winkel von \unit[60]{$^\circ$} innerhalb der definierten Zeitspanne \textit{CpsBlkZeiBlockierungErkannt} überschritten wurde. Zur Unterscheidung zwischen schnellem und langsamem Dekrementieren des Zählers sind die aktuelle Motordrehzahl sowie die Start-Motordrehzahl erforderlich.\\
Die Blockadeauswertung soll sowohl für EC- als auch DC-Motoren funktionsfähig sein. Hierzu wird ein zusätzlicher Eingang für die Blockadeerkennung bei DC-Motoren benötigt (\textit{Blockade detektiert (NSR)}). Darüber hinaus existiert eine sicherheitsrelevante Blockadeerkennung in der Basissoftware, welche die Auswertung der Blockierfunktion beeinflusst.\\
Für die Blockadeauswertung ist der Motorstatus aus dem Zustandautomaten der Basissoftware erforderlich, da der Blockadezähler nur einmal während des Maschinenbetriebs inkrementiert werden kann. Erst ein Zustandswechsel erlaubt ein erneutes Inkrementieren. Zur Blockadeerkennung bei EC-Motoren ist ebenfalls der Motorstatus notwendig, um zu bestimmen, ob sich die Maschine im \textit{running}- oder \textit{off}-Modus befindet.\\
Als Ausgänge werden der maximale Strom und ein Flag für die aktive Stromreduzierung bereitgestellt. Zusätzlich steht ein weiteres Flag für die Maschinenabschaltung zur Verfügung.

\subsection{Modellierung der Funktion in Matlab}
Dieser Abschnitt behandelt die Modellierung der Blockierfunktion in Matlab. Wie in \autoref{ch_voruntersuchungen} erläutert, wird für jede Teilfunktion ein eigenes Referenced-Subsystem sowie ein individuelles Simulink Data Dictionary erstellt. Im Dictionary werden alle Parameter entsprechend den Anforderungen (\autoref{subsec_blk_analyse}) definiert. Dabei erfolgt die Benennung der Parameter gemäß der in \autoref{subsec_namenskonvention} beschriebenen Namenskonvention. Zusätzlich werden zentrale Signale, wie beispielsweise die Ausgänge der Funktion, ebenfalls im Dictionary hinterlegt. Dies ermöglicht den Zugriff auf die Signale über eine A2L-Datei und damit das Messen der Signale auf der realen Maschine über eine XCP-Schnittstelle. \\
Die vollständige Modellierung der Blockierfunktion ist in \autoref{png_matlab_blk} dargestellt. 
\begin{figure}[H]
	\centering
	\includegraphics[width=1\textwidth]{./Bilder/Modell_Blk.pdf}
	\caption{Matlab Modell: Blockierfunktion}
	\label{png_matlab_blk}
\end{figure}
Auf der linken Seite ist das Variant-Subsystem für die Blockadeerkennung bei EC-Motoren dargestellt. Dieses Subsystem wird ausschließlich bei Maschinen mit EC-Motoren aktiviert. Die Aktivierung oder Deaktivierung der Funktion kann zusätzlich über die maschinenabhängige Option \lstinline{ACTIVATE_BLOCKADE_MONITORING} im Header-File der jeweiligen Maschine erfolgen. \\
Am Ausgang des Subsystems werden zwei Flags bereitgestellt. Ein Flag wird zur Abschaltung der Maschine verwendet, während das andere der Erkennung von Blockaden dient. Wie zuvor beschrieben, kann die Blockadeerkennung über drei verschiedene Mechanismen erfolgen. Zum einen durch die sicherheitsrelevante Blockadeerkennung aus der Basissoftware (\textit{sr\_blockade\_detected}) sowie die Blockadeerkennung für DC-Motoren (\textit{nsr\_blockade\_detected}) und die hier implementierte Blockadeerkennung für EC-Motoren. Diese drei Signale sind logisch über ein ODER-Block verknüpft, wobei das resultierende gemeinsame Signal an die Blockadeauswertung weitergeleitet wird.\\
Innerhalb des Stateflow-Charts zur Blockadeauswertung werden alle drei Teilfunktionen innerhalb eines Aufrufes verarbeitet. Dies ist in \autoref{png_matlab_blk_ausw} veranschaulicht.
\begin{figure}[H]
	\centering
	\includegraphics[width=1\textwidth]{./Bilder/Modell_Blk_Auswertung.pdf}
	\caption{Matlab Modell: Auswertung der Blockierfunktion}
	\label{png_matlab_blk_ausw}
\end{figure}
Die obere Teilfunktion ist für das Dekrementieren des Zählers verantwortlich. Dabei wird zwischen einem langsamen und einem schnellen Herunterzählen unterschieden, abhängig davon, ob die aktuelle Motordrehzahl die Startdrehzahl überschreitet. Sobald eine der beiden Zeiten abgelaufen ist, wird der Zähler um eins dekrementiert und die Zeiten zurückgesetzt.\\
Die mittlere Teilfunktion übernimmt das Inkrementieren des Blockadezählers. Wird, wie zuvor beschrieben, über eine der drei Erkennungsfunktionen eine Blockade identifiziert, wird der Zähler (sofern das Maximum noch nicht erreicht ist) um eins erhöht. Ein erneutes Inkrementieren ist erst möglich, wenn der Motorstatus den Standby-Modus erreicht hat. Dieser Zustand tritt ein, sobald die Maschine für eine kurze Zeit ohne Gasbetätigung im Stillstand verbleibt.\\
Die untere Teilfunktion steuert das Aktivieren und Deaktivieren der Stromreduktion basierend auf dem aktuellen Wert des Blockadezählers. Ist die Stromreduktion nicht aktiv, wird der maximale Strom mit dem Wert \textit{BASIS\_CURRENT} ausgegeben. Im Falle einer aktiven Stromreduktion wird der hierfür spezifizierte Parameter der Blockierfunktion verwendet.

%\begin{figure}[H]
%	\centering
%	\includegraphics[width=1\textwidth]{./Bilder/Modell_Blk_ErkEC.pdf}
%	\caption{Matlab Modell: Blockadeerkennung von EC-Motoren}
%	\label{png_matlab_blk_erk}
%\end{figure}